\section*{The Propers: Themes and Movement}

\subsection*{Blessing Before Mastery\enspace\properrefs{Int., Coll., Ep., Grad., All., Gosp., Off., Sec., Comm., Postcomm.}}

Trinity Sunday does not begin by inviting the worshipper to construct a model of God. It begins from Tobit's command to bless God for mercy shown, gives that command an explicitly Trinitarian liturgical form, and then places Paul's astonishment, Daniel's praise, and Christ's baptismal mandate inside that prior thanksgiving. The intellectual order is therefore also a spiritual order: gift, wonder, confession, offering, and mission.

\begin{center}
\begin{tikzpicture}[x=3.55cm]
\node[guide node,text width=2.65cm] (gift) at (0,0) {\textbf{Mercy received}\\Tobit: bless and confess};
\node[guide node,text width=2.65cm] (wonder) at (1,0) {\textbf{Wonder disciplined}\\Romans: depth and doxology};
\node[guide emphasis,text width=2.65cm] (name) at (2,0) {\textbf{One Name confessed}\\Matthew: baptism and teaching};
\node[guide node,text width=2.65cm] (offering) at (3,0) {\textbf{Offerers perfected}\\Secret: victim and self-gift};
\foreach \x/\y in {gift/wonder,wonder/name,name/offering}
  \draw[guide arrow] (\x) -- (\y);
\end{tikzpicture}
\end{center}

\noindent Received mercy, disciplined wonder, the confessed Name, and the perfected offerers make one reading of the formulary from entrance to offering. The cited patristic witnesses directly explain particular biblical texts or illuminate the doctrine of sacrifice; none of them expounds the Roman proper as a whole.

\begin{massbox}{The Collect supplies the rule of reading}
The prayer asks for true faith to \emph{acknowledge} the glory of the eternal Trinity, to \emph{adore} the Unity in the power of majesty, and to be guarded by steadfastness in that faith. Knowledge is ordered to worship, and worship asks for perseverance. Mystery exceeds comprehension without becoming either contradiction or silence.
\end{massbox}

\sectionguard
\subsection*{Romans 11: Wonder Inside the Economy of Mercy\enspace\properrefs{Ep.}}

Paul's cry, ``O the depth,'' concludes Romans 9--11 rather than interrupting the letter with an abstract reflection. He has traced Israel's privileges and stumbling, the Gentiles' unexpected grafting, the warning against Gentile boasting, and the divine purpose to have mercy upon all. The judgments are unsearchable precisely because the apostle has followed a real saving economy farther than creaturely foresight can govern it.

\begin{sourcecard}{St. John Chrysostom --- \emph{Homily 19 on Romans}}
Commenting directly on Romans 11:33--36, Chrysostom says that Paul knows there is a depth but does not claim to know its measure. He follows the reversals of Jew and Gentile, marvels that apparent opposites are governed by one wisdom, and ends where Paul ends: God is no debtor, because every creature first receives from him. Wonder is not a retreat from exegesis; it is exegesis brought to the limit set by divine freedom and beneficence.
\end{sourcecard}

\begin{threestable}{Paul's phrase}{Chrysostom's direct emphasis}{Work within this feast}
\latin{O altitudo} & The apostle can recognize the reality of the depth without measuring it. & Confession is neither mastery nor agnosticism: what is revealed is truly known and inexhaustible.\\
Unsearchable judgments and ways & The surprising dispensations toward Israel and the nations display one wise power through reversals. & The Gospel's mission to all nations is heard within salvation history, not as a detached slogan of universality.\\
Who first gave to him? & God is always the first benefactor and never a creature's debtor. & Tobit's remembered mercy and the altar's offering both begin in prior gift.\\
From, through, and unto him & Paul's inquiry resolves in glory rather than possession. & Doxology is the proper terminus of theological attention and the beginning of obedient worship.\\
\end{threestable}

\begin{studybox}{Romans 11:36 in its own grammar}
Romans 11:36 speaks doxologically of God in its immediate grammar. The verse is fittingly contemplated on Trinity Sunday, but its grammar does not make \latin{ex ipso, et per ipsum, et in ipso} three labels for Father, Son, and Holy Ghost. The Church's Trinitarian confession is governed explicitly by the baptismal Name in Matthew 28 and by the wider rule of faith.
\end{studybox}

\sectionguard
\subsection*{Psalm 8 and Daniel: Praise Learns the Name\enspace\properrefs{Int., Grad., All.}}

The Introit joins Tobit's thanksgiving to ``O Lord our Lord, how admirable is thy name in the whole earth'' (Ps. 8:2). Augustine expounds that verse directly, though not its later placement in this Mass.

\begin{sourcecard}{St. Augustine --- \emph{Exposition on Psalm 8}, 4--8}
Augustine reads the admirable Name in all the earth in the light of Christ's exaltation, whose glory is above the heavens. The psalm then places praise in the mouths of the little ones before mature comprehension. For Augustine, faith and praise are not substitutes for knowledge; they are the creature's truthful beginning before a glory that exceeds every verbal account.
\end{sourcecard}

The Gradual and Alleluia continue Daniel 3's blessing from abyss to throne and from the God of the fathers to everlasting praise. The Gradual does not solve Paul's unsearchable depth. It changes the observer's position: what the creature cannot sound is entirely beheld by the Lord enthroned above the Cherubim. The Alleluia then makes inherited praise ready for Christ's command to carry the one Name to the nations.

\begin{threestable}{Chant movement}{Textual action}{Read across the formulary}
Introit & Mercy is remembered and the admirable Name blessed throughout the earth. & Grace precedes inquiry; the earth-wide Name anticipates the Gospel's nations.\\
Gradual & God beholds the abysses and reigns from the blessed throne of glory. & Paul's ``depth'' is placed under divine sight rather than human control.\\
Alleluia & The God of the fathers is blessed for ever. & Israel's received praise opens toward the mission to all nations.\\
Offertory and Communion & Father, Son, and Holy Ghost are blessed; God is confessed before all the living. & Praise received from Scripture is returned at the altar and reopened toward public witness.\\
\end{threestable}

\sectionguard
\subsection*{The One Name Becomes Observance and Gift\enspace\properrefs{Gosp., Off., Sec., Comm., Postcomm.}}

Matthew's singular ``Name'' names Father, Son, and Holy Ghost without confusing the Persons or dividing the majesty. Hilary, Basil, and Gregory Nazianzen read the command as a rule of faith and initiation: disciple the nations through baptism and teaching under the risen Son's abiding presence. Confession must become observance.

The remaining propers give obedience sacrificial shape. The Offertory adapts Tobit's praise; the Secret asks God, through his holy Name, to sanctify the victim and, through it, to perfect the offerers as enduring gift. Augustine sets the order: visible sacrifice signifies the redeemed city's offering through Christ the High Priest. Communion returns to public confession; the Postcommunion asks that sacrament and Trinitarian faith advance salvation in body and soul.

\begin{studybox}{The revealed Name and the one sacrifice}
The Trinity is confessed from Christ's revealed Name, not reverse-engineered from Paul's prepositions. The offerers become no second saving victim; grace conforms them to Christ's one priestly self-offering.
\end{studybox}
